\documentclass[10pt,a4paper]{amsart}
\usepackage[margin=19mm]{geometry}
\usepackage{amsmath,amssymb,mathtools}
\usepackage[scr=rsfs,cal=euler]{mathalfa}
\usepackage{xfrac}
\usepackage[unicode,hidelinks,bookmarks=false]{hyperref}
\newcommand{\ie}{i.e.\ }
\newcommand{\cf}{cf.\ }
\newcommand{\resp}{respectively\ }
\newcommand{\Subsection}{\subsection}
\providecommand{\nobreakdash}{\nobreak\hbox{-}}
\setlength{\emergencystretch}{2em}
\multlinegap0pt
\usepackage{enumerate}
\usepackage{amssymb}
\usepackage[shortlabels]{enumitem}
\usepackage{tikz}
\usepackage{xcolor}
\newtheorem{lemma}{Lemma}
\newtheorem{thm}{Theorem}
\title{The Manickam-Mikl\'os-Singhi Property in Graphs and Hypergraphs}
\author{Adam D\v{z}avoronok}
\date{Source: arXiv:2605.22708v2 (CC BY 4.0)}
\begin{document}
\maketitle

\noindent\emph{Source and licence.} The Manickam-Mikl\'os-Singhi Property in Graphs and Hypergraphs. Version arXiv:2605.22708v2 (CC BY 4.0), distributed by arXiv under the Creative Commons Attribution 4.0 International licence, http://creativecommons.org/licenses/by/4.0/, as recorded in the arXiv licence metadata for this exact version and shown on the abstract page https://arxiv.org/abs/2605.22708v2. Full licence notice: \texttt{licenses/arxiv-2605.22708-CC-BY-4.0.txt}.\par\smallskip

\noindent\textit{Editorial scope.} Counted unit: the article's thm thm:circulant (source0.tex lines 294--298). The article's complete original proof of it is delivered with it (source0.tex lines 299--336, one \texttt{proof} environment of 38 lines). No mathematics is written, repaired or re-derived: the statement and the proof are the article's own lines, in the article's own notation, and no retained line is substituted or re-flowed.  Retained uncounted as the source-local context the proof needs: lines 276--279.  Nothing was omitted from any retained span, so the package carries no omission record.  No retained line contains a citation, so the wrapper carries no bibliography and no bibliography entry.  No retained line contains a figure environment, an image-inclusion command or drawing code, so no image is delivered and the package declares no asset dependency.  Editorial formatting only, in addition: an AMS \texttt{amsart} preamble that restates the article's own declarations and macros used by the retained lines, namely the environments \texttt{lemma}, \texttt{thm} on separate counters, together with the standard packages those lines need.  The retained environments print in this document as Lemma 1, Theorem 1 in order of appearance, whereas the article prints the same items with the numbering of its own full document.  The complete native source of the article as distributed in the arXiv e-print for this exact version is bundled unchanged as \texttt{sources/201154/source0.tex}.
 \begin{lemma}\label{lem:unique-cycle}
  Suppose $f:V(C_{2m+1})\rightarrow \mathbb{R}$ is such that $\sum_v f(v)\ge0$ and there is exactly one edge $uv\in E(C_{2m+1})$ satisfying $f(u)+f(v)\ge0$.
 Then $f(u),f(v)\ge0$ and furthermore, for each $xy\in E(C_{2m+1})\setminus\{uv\}$, the values $f(x)$ and $f(y)$ cannot be both nonnegative or both negative.
 \end{lemma}

 \begin{thm}\label{thm:circulant}
A coprime circulant graph $C_{2m+1}^{t_1,t_2,\dots,t_k}$ has the MMS property if and only if for each $i\in [k]$ the set
\[\{ |t_i^{-1}t_1|, |t_i^{-1}t_2|,\dots, |t_i^{-1}t_k|\}\pmod{2m+1}\not\subseteq\{1,3\},\]
where $|\ell|=\min(\ell,2m+1-\ell)$ for $\ell\in\mathbb{Z}_{2m+1}$.
 \end{thm}

\begin{proof}

 Let $C_{2m+1}^{t_1,t_2\dots ,t_k}$ be a coprime circulant graph satisfying the described condition. We show it has the MMS property.
 First, we observe that each cycle gives us at least one nonnegative edge.
 On the other hand, if each cycle always has $2$ nonnegative edges, then the coprime circulant graph has the MMS property.
 Hence, we can assume that there exists a function and a cycle with just $1$ nonnegative edge.
 Without loss of generality, assume that this is the cycle generated by $t_1$.
 Consider the relabelling
\[
\varphi : \mathbb{Z}_{2m+1} \to \mathbb{Z}_{2m+1}, \qquad \varphi(x)=t_1^{-1}x.
\]
Since $\gcd(t_1,2m+1)=1$, the element $t_1$ is invertible modulo $2m+1$, so
$\varphi$ is a permutation of the vertex set.
 If $\{x,y\}$ is an edge in the cycle
generated by $t_i$, then
\[
x-y \equiv \pm t_i \pmod{2m+1},
\]
and therefore
\[
\varphi(x)-\varphi(y)=t_1^{-1}(x-y)\equiv \pm t_1^{-1}t_i \pmod{2m+1}.
 \]
Thus, after relabelling, we obtain the graph
\(
C_{2m+1}^{1,t_1^{-1}t_2,\dots,t_1^{-1}t_k},
\)
and the cycle with exactly one nonnegative edge is now the cycle generated by $1$.
 Transporting the weight function along $\varphi$ preserves the number of
nonnegative edges in each cycle.
 We may therefore assume that the cycle generated by $1$ has exactly one
nonnegative edge.
 We fix an order of vertices, such that the vertices $u,v$ from Lemma~\ref{lem:unique-cycle} are $1$ and $2m+1$, then by Lemma~\ref{lem:unique-cycle} the nonnegative vertices are exactly the odd numbers.
 We can now consider the other cycle generated by some $t_1^{-1}t_i$.
 If $|t_1^{-1}t_i|=3$, then a direct verification shows that this cycle has exactly $2$ nonnegative edges given by $\{1,2m-1\}$ and $\{2m+1,3\}$.
 In the other cases, where $|t_1^{-1}t_i|\ge5$ or even, we have at least $3$ nonnegative edges as seen by $\{i,i+|t_1^{-1}t_i|\}$ if $|t_1^{-1}t_i|$ is even or $\{2m+2-i,2m+2-i+|t_1^{-1}t_i|\}$ in the odd case, where $i\in\{1,3,5\}$.
 So overall, for the described coprime circulant graphs, there is at least one cycle with three nonnegative edges and exactly one cycle with one nonnegative edge.
 This proves one direction. The converse follows from the same argument.
\end{proof}

\end{document}
